\documentclass{article}
\usepackage{amsmath}
\usepackage{amssymb}
\usepackage{amsthm}
\usepackage[T1]{fontenc}


\newtheorem*{conjecture}{Conjecture}
\newtheorem*{theorem}{Theorem}
\newtheorem{lemma}{Lemma}
\newtheorem*{lemma*}{Lemma}
\newtheorem*{proposition}{Proposition}


\makeatletter
\newcommand*{\rom}[1]{\expandafter\@slowromancap\romannumeral #1@}
\makeatother



\begin{document}

\section*{Problem 1}

We use multiplication rule to get the number of ways we can order the coffee.
We have:

\begin{equation*}
    \mbox{Choosing 4 types of coffee} - 4 \ \mbox{ways}
\end{equation*}
\begin{equation*}
    \mbox{Choosing 3 different cups} - 3 \ \mbox{ways}
\end{equation*}
\begin{equation*}
    \mbox{Choosing a combination of additions (cream etc.)} - 2^{3}
\end{equation*}
We notice that we also include the possibility of choosing no additions.
Then by the multiplication rule we have:
\begin{equation*}
    4 \cdot 3 \cdot 2^{3} = 12 \cdot 8 = 96
\end{equation*}
ways of ordering the coffee.

\section*{Problem 2}
We have \(n=12\) chairs and we need to assign \(k=8\) of them to \(k=8\) committee
members. We recognize sampling without replacement (each time we choose a chair
we have one less committee member to assign them a chair). We have then:
\begin{equation*}
    \frac{n!}{(n-k)!} = \frac{12!}{(12-8)!} = \frac{12!}{4!} = 12 \cdot 11 \cdot \dots 5 = 19958400
\end{equation*}
ways for the committee to sit in all available chairs.


\subsection*{Appendix - a different way to think about ways}
We can also first find the total number of ways we can choose
"different" chairs from \(n=12\) chairs for \(k=8\) committee members:
\begin{equation*}
    \binom{12}{8} = \frac{12!}{8!(12 - 4)!}
\end{equation*}
and knowing that it does not preserve the order, we multiply that
by \(8!\) possible ways of rearranging the committee members (there is 
now \(8\) members sitting in \(8\) positions) so the total number of ways is:
\begin{equation*}
    \binom{12}{8} \cdot 8! = \frac{12!}{8!(12 - 8)!} \cdot 8! = \frac{12!}{4!} = 19958400
\end{equation*}
in agreement with the first way.
\end{document}
